\begin{figure}[H]
\centering
\includegraphics[width=0.95\linewidth]{figures/04_orbital_mapper_interpretation.pdf}
\caption{Interpretacion del espacio orbital.}
\end{figure}

La Figura 6 presenta el caso mas convincente del reporte. El grafo orbital combina un nucleo dominante con varias ramas y pequenos componentes satelite. Esa mezcla es exactamente lo que vuelve interesante a Mapper en este contexto: no solo hay densidad de objetos en una region central, sino tambien separaciones locales que podrian corresponder a regimes orbitales distintos, por ejemplo sistemas compactos de periodo corto frente a poblaciones mas dispersas de periodo largo. Como ademas la imputacion media en este espacio es muy baja, estas separaciones merecen seguimiento cientifico detallado.

\begin{figure}[H]
\centering
\includegraphics[width=0.95\linewidth]{figures/05_joint_mapper_interpretation.pdf}
\caption{Interpretacion del espacio conjunto.}
\end{figure}

La Figura 7 sugiere una lectura distinta para el espacio conjunto. Aqui la topologia sigue siendo rica, pero visualmente el grafo esta mas concentrado en un componente central compacto, con menos periferia interpretable que en orbital. Esto es consistente con una mezcla de senales orbitales y fisicas que aumenta la densidad del centro del espacio embebido. El resultado sigue siendo util porque conecta masa, radio y orbitalidad, pero tambien parece mas sensible al lens y a la eleccion de variables.

\begin{figure}[H]
\centering
\includegraphics[width=0.95\linewidth]{figures/09_thermal_caution.pdf}
\caption{Complejidad alta pero confianza reducida en el espacio termico.}
\end{figure}

La Figura 8 resume la principal advertencia del trabajo. El espacio termico genera un grafo grande y visualmente muy articulado, pero la coloracion por imputacion muestra que buena parte del componente central y de la periferia se apoya en fracciones elevadas de datos completados. En otras palabras, el grafo termico es interesante como detector de hipotesis, pero aun no como evidencia fuerte de poblaciones termicas bien separadas.
